\documentclass[a4paper,12pt]{article}
\usepackage[utf8]{inputenc}
\usepackage[russian]{babel}
\usepackage{amsmath, amssymb}
\usepackage{algorithm}
\usepackage{algpseudocode}
\usepackage{geometry}
\geometry{a4paper, margin=2cm}

\title{Алгоритм топологической пространственной фильтрации (fit\_filters)}
\author{}
\date{}

\begin{document}

\maketitle

\begin{abstract}
В данном документе представлено подробное математическое и алгоритмическое описание метода \texttt{fit\_filters}, предназначенного для оптимизации топологических пространственных фильтров. Алгоритм обучается находить пространственные веса, которые проецируют ковариационные матрицы таким образом, чтобы логарифм мощности выходного сигнала в пространстве меньшей размерности аппроксимировал топологическую структуру заданного многообразия, следуя принципам Uniform Manifold Approximation and Projection (UMAP).
\end{abstract}

\section{Введение}
Функция \texttt{fit\_filters} реализует процедуру обучения набора (батча) топологических пространственных фильтров, используя принципы дифференцируемого UMAP.
Входными данными являются ковариационные матрицы $\mathbf{C}_i \in \mathbb{R}^{M \times M}$ (где $M$ --- количество каналов, $i \in \{1, \dots, N_{epochs}\}$), а целевой топологией выступает граф соседства, построенный либо по исходным признакам $\mathbf{T}$, либо по заданной матрице расстояний $\mathbf{D}$.

Алгоритм одновременно обучает $K$ независимых инициализаций (restarts) фильтров размерности $N_{dim}$, чтобы избежать локальных минимумов.

\section{Математическая постановка задачи}

\subsection{Топологический пространственный фильтр}
Для каждого restart $k \in \{1, \dots, K\}$ и каждой размерности $d \in \{1, \dots, N_{dim}\}$ пространственный фильтр определяется вектором весов $\mathbf{w}_{k,d} \in \mathbb{R}^{M}$ и скалярным параметром масштаба $s_{k,d} > 0$.

Проекция $i$-й ковариационной матрицы $\mathbf{C}_i$ в одномерное пространство для фильтра $\mathbf{w}_{k,d}$ (мощность сигнала после фильтрации) вычисляется как квадратичная форма:
\begin{equation}
    P_{k,i,d} = \mathbf{w}_{k,d}^T \mathbf{C}_i \mathbf{w}_{k,d}
\end{equation}

Координата вложения $y_{k,i,d}$ получается путем логарифмирования мощности и применения обучаемого масштаба:
\begin{equation}
    y_{k,i,d} = s_{k,d} \cdot \log(P_{k,i,d} + \epsilon)
\end{equation}
где $\epsilon = 10^{-8}$ --- малая константа для вычислительной стабильности, а $s_{k,d} = \exp(\log s_{k,d})$, что гарантирует его положительность.

Векторная координата вложения для $i$-й эпохи и $k$-го restart:
\begin{equation}
    \mathbf{y}_{k,i} = [y_{k,i,1}, \dots, y_{k,i,N_{dim}}]^T \in \mathbb{R}^{N_{dim}}
\end{equation}

\subsection{Целевая топология (UMAP Graph)}
На основе входных признаков или матрицы расстояний строится нечеткий симплициальный граф (fuzzy simplicial set). Матрица смежности графа $V \in \mathbb{R}^{N_{epochs} \times N_{epochs}}$ содержит вероятности $v_{ij} \in [0, 1]$ наличия ребра между эпохами $i$ и $j$.

\subsection{Функция потерь (UMAP Cross-Entropy)}
Для каждого restart $k$ рассчитывается функция потерь, представляющая собой перекрестную энтропию (cross-entropy) между исходным графом $v_{ij}$ и графом, индуцированным попарными расстояниями в пространстве вложений $\mathbf{y}_{k,i}$.

Квадрат евклидова расстояния в пространстве вложений:
\begin{equation}
    d_{k,ij}^2 = \|\mathbf{y}_{k,i} - \mathbf{y}_{k,j}\|_2^2
\end{equation}

Вероятность связи в пространстве вложений (согласно UMAP):
\begin{equation}
    w_{k,ij} = \frac{1}{1 + a (d_{k,ij}^2)^b}
\end{equation}
где $a$ и $b$ --- гиперпараметры UMAP, определяющие форму кривой затухания.

Функция потерь для $k$-го restart:
\begin{equation}
    \mathcal{L}_k = \sum_{i \neq j} \left( v_{ij} \log \left( \frac{v_{ij}}{w_{k,ij}} \right) + (1 - v_{ij}) \log \left( \frac{1 - v_{ij}}{1 - w_{k,ij}} \right) \right)
\end{equation}
Общая функция потерь для батча является средним по всем $K$ restarts: $\mathcal{L}_{UMAP} = \frac{1}{K} \sum_{k=1}^K \mathcal{L}_k$.

\subsection{Регуляризация масштаба}
Поскольку фильтры определяются с точностью до масштаба, вводится слабая регуляризация (penalty) на логарифм параметра $s_{k,d}$ (стремление $\log s_{k,d} \to 0$):
\begin{equation}
    \mathcal{R}(\log \mathbf{s}) = \lambda \frac{1}{K} \sum_{k=1}^K \sum_{d=1}^{N_{dim}} (\log s_{k,d})^2
\end{equation}
где $\lambda$ (\texttt{scale\_reg}) --- коэффициент регуляризации.

Итоговая целевая функция:
\begin{equation}
    \mathcal{J} = \mathcal{L}_{UMAP} + \mathcal{R}(\log \mathbf{s})
\end{equation}

\section{Алгоритм оптимизации}

Оптимизация параметров $\mathbf{w}$ и $\log \mathbf{s}$ производится с помощью алгоритма \textbf{AdamW}.

Особое внимание уделяется предотвращению тривиальных решений (когда веса фильтров стремятся к нулю, а масштаб бесконечно возрастает):
\begin{enumerate}
    \item Векторы весов $\mathbf{w}_{k,d}$ жестко нормируются на единичную L2-норму после каждого шага градиентного спуска:
    \begin{equation}
        \mathbf{w}_{k,d} \leftarrow \frac{\mathbf{w}_{k,d}}{\|\mathbf{w}_{k,d}\|_2}
    \end{equation}
    \item Значения $\log s_{k,d}$ ограничиваются в диапазоне $[\log s_{min}, \log s_{max}]$ (отсечение/clamping).
\end{enumerate}

\begin{algorithm}[H]
\caption{Алгоритм \texttt{fit\_filters}}
\begin{algorithmic}[1]
\Require Ковариационные матрицы $\{\mathbf{C}_i\}_{i=1}^{N_{epochs}}$, целевой граф $V$, количество restarts $K$, размерность $N_{dim}$, число эпох $E$, learning rate $\alpha$.
\Ensure Оптимальные фильтры $\mathbf{w}^*$, масштабы $\mathbf{s}^*$.
\State Инициализация $\mathbf{w} \in \mathbb{R}^{K \times N_{dim} \times M}$ (случайные значения из $\mathcal{N}(0, 1)$ или заданные).
\State Инициализация $\log \mathbf{s} \in \mathbb{R}^{K \times N_{dim}}$ (нули или заданные).
\State Настройка оптимизатора AdamW (для $\mathbf{w}$ и $\log \mathbf{s}$) и планировщика ReduceLROnPlateau.
\For{epoch $= 1, 2, \dots, E$}
    \State Вычислить мощность $P_{k,i,d} = \mathbf{w}_{k,d}^T \mathbf{C}_i \mathbf{w}_{k,d}$
    \State Вычислить вложения $\mathbf{y}_{k,i} = \exp(\log s_{k,d}) \cdot \log(P_{k,i,d} + \epsilon)$
    \State Вычислить UMAP-потери $\mathcal{L}_k$ на основе $\mathbf{y}_{k,i}$ и $V$
    \State Вычислить регуляризацию $\mathcal{R} = \lambda (\log \mathbf{s})^2$
    \State $\mathcal{J} = \text{mean}(\mathcal{L}) + \text{mean}(\mathcal{R})$
    \State Обратное распространение ошибки: $\nabla_{\mathbf{w}} \mathcal{J}$, $\nabla_{\log \mathbf{s}} \mathcal{J}$
    \State Обновить веса: AdamW.step()
    \State Проекция весов: $\mathbf{w}_{k,d} \leftarrow \mathbf{w}_{k,d} / \|\mathbf{w}_{k,d}\|_2$
    \State Ограничение масштабов: $\log \mathbf{s} \leftarrow \text{clamp}(\log \mathbf{s}, \min, \max)$
    \State Планировщик LR: шаг на основе $\mathcal{J}$
\EndFor
\State Пост-обработка: для каждого $k$ оценить индивидуальные потери каждой размерности, отсортировать $\mathbf{w}_{k,d}$ внутри каждого restart.
\State Отсортировать $K$ restarts по итоговой функции потерь (от лучшего к худшему).
\State \Return $\mathbf{w}^*$ (лучшие restarts), $\mathbf{s}^*$, истории потерь.
\end{algorithmic}
\end{algorithm}

\section{Заключение}
Описанный алгоритм позволяет находить проекции ковариационных матриц сигналов, в которых логарифм мощности сохраняет нелинейную многообразную структуру исходного пространства признаков или матрицы расстояний. Использование нескольких независимых стартов (restarts) повышает робастность нахождения глобального или хорошего локального минимума сложной целевой функции UMAP.

\end{document}